\chapter{Datos y mediciones a lo largo del año}\label{Capitulo_3}
\chapterquote{}{}
\section{Introducción}\label{III:introduccion}

Desde hace mucho tiempo se sabe que el comportamiento de las distintas especies animales varía con la sucesión de las estaciones. 
Este conocimiento puede ser tan fundamental como comprender las migraciones de las presas, en épocas en las que el ser humano vivía 
de la caza, o como saber que los osos hibernan durante los meses de bajas temperaturas. En el caso de la tortuga terrestre argentina, 
esta variación estacional se manifiesta de forma marcada: durante el invierno los individuos reducen su actividad y entran en brumación, 
resguardándose en refugios excavados en el suelo.

Si bien estos conocimientos son más propios de las ciencias naturales, la biología y la ecología, la cuantificación de los cambios estacionales
en el comportamiento de los animales es un tema que puede abordarse desde la perspectiva de la física o del análisis de datos. Por ello, en este
capítulo vamos a presentar los datos obtenidos a lo largo de un año con las mediciones de GPS, nombradas en el capítulo \ref{Capitulo_1}, 
para posteriormente ofrecer métricas cuantitativas que caractericen la evolución del comportamiento de esta especie. El objetivo de este capítulo, 
será intentar observar el comportamiento de brumación a partir de las métricas evaluadas, como también brindar algún aporte que contribuyan al 
entendimiento de los métodos de la selección de refugios nocturnos.

\section{Datos de GPS}
Tal como fue contado en el primer capítulo de la tesis, los datos de GPS corresponden a la experiencia de campo durante la realización tesis de la 
tesis de maestría M. Madile \cite{madile2023} en 2022-2023. A estos datos se les tuvo que aplicar un filtrado para poder eliminar algunos saltos entre 
mediciones continuas de ubicaciones, las cuales serían imposible de realizar por tortugas terrestre que tienen un desplazamiento lento. 

La principal técnica que se utilizó para descartar mediciones inconsistentes, fue a partir de poner una cota máxima ($v_{\max}$) a la velocidad promedio entre mediciones.
Por ello se tomaron dos mediciones consecutivas con una diferencia temporal $\delta t$ y a una distancia $\delta x$, entonces si $\frac{\delta x}{\delta t} \geq v_{\max}$ 
se descartaba la última de las dos mediciones, dado que al hacer este proceso en orden cronológico, la primera ya fue validada. Fijamos este umbral en 
$v_{\max}=0.1\,\mathrm{m/s}$, lo que equivale a un desplazamiento máximo de aproximadamente $90\,\mathrm{m}$ entre mediciones separadas por 15 minutos.
La segunda técnica utilizada consistió en descartar mediciones por fuera del rango de latitudes correspondiente a la zona de estudio, $-40.60^\circ$ a 
$-40.55^\circ$, descartando en torno al $ 8.6\% $ de los datos (Tabla \ref{tab:filtrado}).


\begin{table}[ht]
    \centering
    \caption{Mediciones de GPS descartadas por el filtrado, para cada individuo.\label{tab:filtrado}}
    \begin{tabular}{lrrrr}
    \hline
    Individuo & Mediciones & Descartadas & \% & Conservadas \\
    \hline
    T10 & 12\,321 & 2\,405 & 19.5 & 9\,916 \\
    T12 & 13\,712 & 1\,753 & 12.8 & 11\,959 \\
    T30 & 10\,699 & 56 & 0.5 & 10\,643 \\
    T54 & 15\,094 & 1\,954 & 13.0 & 13\,140 \\
    T6  & 15\,157 & 662 & 4.4 & 14\,495 \\
    T79 & 13\,233 & 46 & 0.4 & 13\,187 \\
    \hline
    Total & 80\,216 & 6\,876 & 8.6 & 73\,340 \\
    \hline
    \end{tabular}
\end{table}

La proporción de mediciones descartadas varía considerablemente entre individuos (Tabla \ref{tab:filtrado}): mientras que en T30 y T79 el ruido es 
prácticamente inexistente ($ < 0.6\%$), en T10 alcanza casi el $20\%$. Esta disparidad sugiere que el ruido depende principalmente del dispositivo 
y no de la conducta del animal, lo cual permite descartar sin preocuparnos de estar sobreeliminando datos.


\section{Detección de refugios}\label{G:refugios}

También realizamos análisis sobre los refugios utilizados por los individuos, por lo que fue necesario definirlos a partir de las ubicaciones proporcionadas
por los datos de GPS filtrados, como nombramos previamente. Dicho trabajo no fue trivial de realizar, ya que como mencionamos en el Capítulo \ref{Capitulo_1}, 
hubo un largo período en el año donde no se midieron las posiciones durante las noches. Se sabe que las tortugas suelen pasar la noche en un mismo refugio 
sin moverse, por ello optamos por descartar para la detección de refugios las mediciones nocturnas (20:00 a 6:00) cuya dispersión media sea $\sigma \geq 10\,\mathrm{m}$.
Esta cota de $10\,\mathrm{m}$ es del orden del error de posicionamiento del GPS y permite distinguir un animal quieto de una medición dominada por ruido.
En caso de que haya al menos 3 datos nocturnos cuya dispersión no supere el umbral, se registró la posición media y se la clasificó como un refugio. 
Para las noches en las que no se dispone de ninguna medición en esa ventana, utilizamos como alternativa las mediciones diurnas entre las 10 y las 18 h, 
aplicando el mismo criterio. Finalmente, agrupamos las posiciones de los refugios de cada individuo mediante DBSCAN con un radio de $10\,\mathrm{m}$ y 
un mínimo de una medición, de modo que noches con posiciones a menos de esa distancia se asignan a un mismo refugio.

%%% Local Variables: 
%%% mode: latex
%%% TeX-master: "template"
%%% End: 
